\documentclass[tikz,border=15pt]{standalone}
\usepackage{ctex}      % 提供中文支持
\usepackage{array}
\usepackage{tabularx}
\usepackage{booktabs}  % 优雅的三线表
\usepackage{xcolor}
\usepackage{colortbl}

% 定义表头颜色：一抹沉稳优雅的高级深灰蓝，其余保持极简黑白
\definecolor{headerbg}{RGB}{52, 73, 94} 
\definecolor{headertext}{RGB}{255, 255, 255} 

\begin{document}

% 设置无衬线字体
\sffamily
% 整体行高稍微回调至 1.3，避免过度松散
\renewcommand{\arraystretch}{1.3}

% 第一列加宽至 4.5cm，第二列自动铺满。
\begin{tabularx}{22cm}{>{\bfseries}p{4.5cm} X}
	\rowcolor{headerbg}
	\toprule
	\textcolor{headertext}{\textbf{文件系统}} & \textcolor{headertext}{\textbf{引进机制、优缺点对比}}                                                                           \\
	\midrule

	% 分类 1
	\multicolumn{2}{l}{\textbf{1. Linux 原生文件系统}}                                                                                                                  \\
	\addlinespace[0.5ex]
	ext (已弃用)                             & 首个为 Linux 设计的文件系统。\newline \textbf{优势：}突破 Minix 系统限制。\newline \textbf{劣势：}磁盘碎片化严重，单文件最大仅支持 2GB。                       \\
	\addlinespace[1ex]
	ext2                                  & 引入 inode 表与数据块管理。\newline \textbf{优势：}解决碎片化问题，稳定性高，为早期 Linux 标配。\newline \textbf{劣势：}无日志功能，异常断电后执行 fsck 全盘检查耗时长。      \\
	\addlinespace[1ex]
	ext3                                  & 引入日志功能 (Journaling)。\newline \textbf{优势：}记录数据更改轨迹，提升异常宕机后的数据恢复速度。\newline \textbf{劣势：}日志记录产生额外性能开销，不支持动态磁盘碎片整理。       \\
	\addlinespace[1ex]
	ext4                                  & 引入区段 (Extents) 与延迟分配。\newline \textbf{优势：}提升大文件与海量存储处理能力，为现代多数 Linux 默认标配。\newline \textbf{劣势：}删除文件会清空相关指针，误删后数据恢复困难。 \\
	\addlinespace[1ex]
	XFS (RHEL 默认)                         & 引入分配组 (Allocation Groups)。\newline \textbf{优势：}64位企业级日志文件系统，并发读写性能优异，适合处理超大文件。\newline \textbf{劣势：}文件系统体积仅支持扩容，不支持缩减。 \\
	\addlinespace[1ex]
	ReiserFS (已弃用)                        & 引入 B* 树结构。\newline \textbf{优势：}小文件读写性能优异。\newline \textbf{劣势：}缺乏后续商业支持与社区维护。                                          \\
	\midrule

	% 分类 2
	\multicolumn{2}{l}{\textbf{2. 网络文件系统}}                                                                                                                        \\
	\addlinespace[0.5ex]
	proc                                  & 内存级别的虚拟文件系统，用于动态映射内核状态与进程信息。                                                                                          \\
	\addlinespace[1ex]
	nfs                                   & Unix/Linux 服务器之间挂载共享目录的标准网络协议。                                                                                        \\
	\addlinespace[1ex]
	smb                                   & 用于 Linux 与 Windows 系统之间进行网络文件共享的协议。                                                                                   \\
	\midrule

	% 分类 3
	\multicolumn{2}{l}{\textbf{3. 跨平台文件系统}}                                                                                                                       \\
	\addlinespace[0.5ex]
	ntfs                                  & Windows 系统的默认文件系统，支持大文件与权限控制。(Windows)                                                                                \\
	\addlinespace[1ex]
	vfat                                  & 即 FAT32 文件系统，单文件上限 4GB，跨平台兼容性高。(常用 U 盘格式)                                                                             \\
	\addlinespace[1ex]
	msdos                                 & 早期的 FAT16 文件系统。(早期 DOS / Windows，已极少使用)                                                                               \\
	\addlinespace[1ex]
	iso9660                               & CD/DVD-ROM 光盘标准只读文件系统。(通用光盘格式)                                                                                        \\
	\midrule

	% 分类 4
	\multicolumn{2}{l}{\textbf{4. 早期 Unix 文件系统}}                                                                                                                  \\
	\addlinespace[0.5ex]
	minix                                 & Linux 诞生之初借鉴的教学级文件系统。(Minix，已极少使用)                                                                                    \\
	\addlinespace[1ex]
	jfs / hpfs                            & 分别是 IBM 与 OS/2 的早期高性能文件系统。(AIX / OS/2，已极少使用)                                                                          \\
	\addlinespace[1ex]
	sysv / ufs                            & 早期 System V Unix 与 BSD 的标准文件系统。(老旧 Unix / BSD，已极少使用)                                                                  \\
	\addlinespace[1ex]
	umsdos / ncp                          & 建立在 DOS 上的类 Unix 文件系统 / Novell NetWare 网络文件系统。(遗留系统，已极少使用)                                                            \\
	\bottomrule
\end{tabularx}

\end{document}